\section{Method}\label{sec:method}

The method is a portfolio of three arms: multilevel splitting driven by transient duration (channel B), active learning of basin boundaries (channel A), and a uniform share that bounds the worst case. Prior sets (channel C) enter as a distribution of starting points and candidates. For each point, all arms receive only its trajectory fingerprint and the label produced by the identification protocol (\cref{sec:protocol}); they do not know which label is the target.

\subsection{Reaction coordinate: transient duration}\label{sec:xi}

For an initial point $x_0$, we define
\begin{equation}
\xi(x_0)=\begin{cases}\ln k_{\rm esc}/\ln N_{\rm tot}\in[0,1], & \text{escape at step } k_{\rm esc},\\[1mm] 1+D(x_0), & \text{bounded trajectory},\end{cases}
\label{eq:xi}
\end{equation}
\begin{equation*}
D=\bigl\|(\bar x_{\rm early}-\bar x_{\rm obs},\ {\rm sd}_{\rm early}-{\rm sd}_{\rm obs})\bigr\|/L,
\end{equation*}
where $k_{\rm esc}$ is the integration step at which the trajectory leaves the escape radius, $N_{\rm tot}$ is the number of integration steps, $\bar x$ and ${\rm sd}$ are the coordinate-wise means and standard deviations over the early window (the first steps of the transient segment) and over the observation window, and $L$ is half the largest edge of $B$. A trajectory that settles quickly onto its regime has $D\approx0$; one that lingers near an edge state on a basin boundary has a large $D$ or a large $k_{\rm esc}$. The two branches order all trajectories on one scale: every escaping trajectory ranks below every bounded one, escaping trajectories are ranked by their escape time, and bounded ones by how slowly they settle. The first branch carries the escape-time information of \cref{cor:channels}(b); the second varies inside the nonescaping set $Q_\infty$, where escape times are constant, and it is all that $\xi$ sees on systems whose trajectories are bounded. In rare-event algorithms, the optimal score function is the committor~\cite{Lestang2018TAMS,Lucente2022Committor}, which is unknown here; $\xi$ is a cheap proxy computed from one trajectory without the target label. On uniform samples, $\xi$ separates the points near the target basin from the rest with an area under the ROC curve of 0.88--0.97 on all systems but one, whereas the asymptotic features (the Lyapunov exponent and the maximum norm) reach only 0.5--0.75 on all systems but the synthetic family D1 (section~\ref{app:diagnostics}).

\subsection{Multilevel splitting in the space of initial conditions}\label{sec:sus}

The static counterpart of adaptive multilevel splitting~\cite{CerouGuyader2007} for the map ``initial condition $\mapsto$ fingerprint'' is subset simulation~\cite{AuBeck2001}, which we call splitting. A reference sample of $n_0=50$ uniform points is drawn; the level $\xi_l$ is the $(n_0p_0)$th largest value of $\xi$ with $p_0=0.2$, and these 10 points seed Markov chains. Each chain makes $1/p_0-1=4$ random-walk steps $y=x+\sigma(b-a)\eta$, $\eta\sim\mathcal N(0,I)$, with mirror reflection into the box $B=[a,b]$; a step is accepted if $\xi(y)\ge\xi_l$, so a level costs $n_0(1-p_0)=40$ evaluations, and for a fixed $\sigma$ this is the Metropolis algorithm for the uniform distribution on $\{\xi\ge\xi_l\}$. The step $\sigma$ (initially 0.1) is adapted to the acceptance rate~\cite{Papaioannou2015}, the next level is chosen from all states of the current level, and on stagnation (a relative increase of the level below 2\% or $\sigma<0.005$) for two levels, and after 10 levels, the algorithm restarts from a new reference sample. Restarts protect against chains that have contracted onto a nontarget region with a long transient. Section~\ref{app:ablation} shows the chains on a basin map and compares $\xi$ with the escape step alone and with a multidimensional rarity score.

\subsection{Active learning of basin boundaries}\label{sec:al}

The active-learning arm stores all evaluated points with their labels (escaping trajectories form a class), trains a $k$-nearest-neighbor classifier ($k=10$, inverse-distance weights) in normalized coordinates, and evaluates the candidates with the largest entropy of the predicted class probabilities (uncertainty sampling~\cite{Settles2012}). Half of a pool of 2000 candidates consists of Gaussian perturbations of boundary points of the sample (points with a neighbor of another class), with a scale of half the distance to that neighbor; the other half consists of uniform points or points of a prior set. While the data contain a single class, the arm is idle.

\subsection{Portfolio and worst-case guarantee}\label{sec:portfolio}

In each round, the portfolio evaluates 10 points from the splitting arm, 20 points from the active-learning arm (if the data contain at least two classes), and $\lceil\theta n_{\rm arm}/(1-\theta)\rceil$ uniform points, where $n_{\rm arm}$ is the number of points from the arms and $\theta=0.2$; the points of a round are shuffled, so no arm is favored in the accounting of the first hit. An arm is \emph{autonomous} if the choice of its points depends only on the results of its own evaluations; the uniform and splitting arms are autonomous, and the active-learning arm, which uses all data, is not. The uniform share plays the role of the defensive component of mixture importance sampling~\cite{Hesterberg1995}, and the combination of strategies with different failure modes follows the idea of algorithm portfolios~\cite{Huberman1997,GomesSelman2001}.

\begin{proposition}[portfolio bound]\label{prop:portfolio}
Let $c$ be the largest number of points in a round, $\omega_j$ the guaranteed share of arm $j$ in each round, and $N_{1-\delta}(j)$ the cost of arm $j$ run on its own. Then
\begin{equation}
N_{1-\delta}(\text{portfolio})\le\min\Bigl\{\frac{N^0_{1-\delta}+c}{\theta},\ \min_{j\ \text{autonomous}}\frac{N_{1-\delta}(j)+c}{\omega_j}\Bigr\}.
\label{eq:portfolio}
\end{equation}
\end{proposition}

\begin{proof}
After $N$ evaluations of the portfolio, an autonomous arm $j$ has performed at least $\omega_jN-c$ evaluations of its own, and by autonomy its points are distributed as in a stand-alone run. Hence $\Pr[\tau>N]\le\Pr[\tau_j>\omega_jN-c]\le\delta$ whenever $\omega_jN-c\ge N_{1-\delta}(j)$. For the uniform arm, $N_{1-\delta}(j)=N^0_{1-\delta}$ and $\omega_j=\theta$.
\end{proof}

The proposition is an elementary property of a defensive mixture; its value lies in the construction it justifies. In integrations, the portfolio is at most $1/\theta=5$ times costlier than uniform sampling up to the additive term $c/\theta$, whereas single adaptive methods have no such bound (\cref{sec:failures}). As $p\to0$, its scaling exponent is no larger than that of the best autonomous arm: a constant factor $1/\omega_j$ is lost, but not the order. In our implementation, the share of the splitting arm is $\omega\approx0.26$ when the active-learning arm is active and $\omega\approx0.77$ without it. The bound holds in integrations, not in integrator steps or wall time, because the arms select long transients (\cref{sec:cost}).

The bound applies to every autonomous arm, and in particular to arms that climb the same coordinate in a different way. Stagger-and-step~\cite{SweetNusseYorke2001} keeps one point and proposes $x+\sigma(b-a)v$ with a random unit vector $v$ and a step $\sigma$ log-uniform on $[10^{-3},1]$; the proposal is accepted if $\xi$ increases, and after 30 consecutive rejections the chain restarts from a uniform point. An arm of $n_c$ independent chains, each contributing one point per round, is autonomous.

\begin{corollary}[stagger-and-step arm]\label{cor:sas}
A portfolio that contains an arm of $n_c$ stagger-and-step chains with share $\omega_{\rm sas}$ satisfies, in addition to \eqref{eq:portfolio},
\begin{equation*}
N_{1-\delta}(\text{portfolio})\le\frac{N_{1-\delta}(\text{chains})+c}{\omega_{\rm sas}}.
\end{equation*}
As $p\to0$, its scaling exponent is at most that of the chains, and its cost never exceeds $(N^0_{1-\delta}+c)/\theta$, even where the chains fail.
\end{corollary}

The corollary makes the choice between splitting and stagger-and-step a matter of constants rather than of risk: whichever climbs the coordinate better, the portfolio inherits its order, and the uniform share bounds the loss where the coordinate is deceptive. We test it in \cref{sec:sas}.

\textbf{Two-stage evaluation.} Most integrations are spent on trajectories that settle early onto a known regime. For the cost measurements in steps and wall time (\cref{sec:cost}), every initial condition, for all strategies including uniform sampling, is evaluated in two stages: a short first stage stops trajectories that are identified with confidence as known nontarget regimes, and all others continue to the full horizon, so that their fingerprints coincide bit for bit with those of a single full integration. The rule, its checks, and the implementation are given in section~\ref{sec:twostage}. The gains in integrations of \cref{sec:coupled,sec:scaling_p,sec:rare_ab,sec:sas} are measured with full integration, except those of the coupled Chua circuits, which were run only in the two-stage configuration.
